\documentclass{article}
\begin{document}
\section*{Boolean reduction cost audit, 2026-10-10}
P versus NP remains OPEN. This replicates and narrows earlier BVE portfolio work (ARC-STRENGTHEN5), not a new unrestricted theorem.
\paragraph{Literature.} Een and Biere, Effective Preprocessing in SAT through Variable and Clause Elimination (SAT 2005), prove satisfiability-preserving variable elimination with resolution and practical bounded elimination. Biere, Jarvisalo and Kiesl, Preprocessing in SAT Solving (2021), survey blocked-clause and proof-producing preprocessing.
\paragraph{Exact cost.} Eliminating variable x by Davis-Putnam resolution forms pairs between P positive and Q negative occurrence clauses. The naive generation count is PQ, and the distinct non-tautological resolvents can also be PQ. For disjoint auxiliary variables, clauses (x or a_i) for i=1..p and (not x or not b_j) for j=1..p generate p squared distinct resolvents from 2p clauses. Thus non-growth BVE rejects the pivot for p>2. This is a restricted-rule counterexample, not SAT hardness.
\paragraph{Test.} Seed 20261010, 270 random 3-CNF instances n=3..8, plus four explicit blowup controls p=2..5. Exhaustive truth-table SAT checks after every attempted elimination passed, 274 cases total, 0.103 seconds. Gate optimality, worst-case asymptotics and all-algorithm transfer were not tested.
\paragraph{Open target.} Construct an exact, source-first compact representation of resolvent cross products with charged bit access, discovery, and recursive semantics. Test against solver bypass and the all-input fixed-power shrinkage requirement. No universal SAT successor is established.
\end{document}
